\documentclass[UTF8,fontset=windows,12pt]{ctexart}
%\setcounter{tocdepth}{1}
%\setcounter{secnumdepth}{4}
%\usepackage{ctex} %若要输入中文请取消注释或者把article改成ctexart
% This first part of the file is called the PREAMBLE. It includes
% customizations and command definitions. The preamble is everything
% between \documentclass and \begin{document}.
%\usepackage[OT1]{fontenc}
%\usepackage{fontspec}
%\setmainfont{Times New Roman}
%\setCJKmainfont{Microsoft YaHei}
\usepackage[left=1.91cm,right=1.91cm,top=2.54cm,bottom=2.54cm]{geometry}
\usepackage{amsfonts}
\usepackage{amsmath}
\usepackage{amssymb}
\usepackage[upint,lcgreekalpha]{stix}
\usepackage[type1]{garamondlibre}
%\usepackage{esint}
\usepackage{pgfplots}
\pgfplotsset{compat=1.18}
\usepackage{tikz}
\usetikzlibrary{arrows}
\usepackage{extarrows}
\usepackage[only,llbracket,rrbracket]{stmaryrd}
\usepackage{titlesec}
\usepackage{bm}
\usepackage{graphicx}
\usepackage{appendix}
\usepackage{tikz-cd}
\usepackage{tikz-3dplot}
\usetikzlibrary{arrows.meta, decorations.pathreplacing, calc, patterns}
\usepackage{float}
\usepackage{mathtools}
\usepackage{amsthm}
\usepackage{verbatim}              % better theorem environments
\usepackage{setspace}
\usepackage{enumitem}
\setenumerate[1]{itemsep=0pt,partopsep=0pt,parsep=\parskip,topsep=5pt}
\setitemize[1]{itemsep=0pt,partopsep=0pt,parsep=\parskip,topsep=5pt}
\setdescription{itemsep=0pt,partopsep=0pt,parsep=\parskip,topsep=5pt}
\usepackage{xcolor}
\usepackage[colorlinks,linkcolor=black,anchorcolor=black,citecolor=black]{hyperref}
%\usepackage{apacite}
\usepackage[numbers]{natbib}
%\usepackage{showkeys}
\usepackage{cases}
\usepackage{fancyhdr}
\usepackage[scaled=0.92]{helvet}	% set Helvetica as the sans-serif font
%\usepackage{upgreek}
%\renewcommand{\rmdefault}{ptm}		% set Times as the default text font
%\usepackage{newtxtext}
\bibliographystyle{plain}
% If AMS-LaTeX is used, it can be loaded before or after mtpro2		
% The following loads metro and defines some common MTPro options [2, 4]
%\usepackage[lite,subscriptcorrection,nofontinfo]{mtpro2}
%\pdfmapfile{+mtpro2.map}

%\usepackage{txfonts}

% various theorems, numbered by section

\makeatletter
\renewenvironment{proof}[1][\proofname]{\par
  \pushQED{\qed}%
  \normalfont \topsep6\p@\@plus6\p@\relax
  \trivlist
  \item[\hskip\labelsep
        \bfseries % 关键修改：斜体改为加粗
        #1\@addpunct{.}]\ignorespaces
}{%
  \popQED\endtrivlist\@endpefalse
}
\makeatother
\newtheoremstyle{kaiti-theorem}   % 样式名称
  {3pt}                           % 上方间距
  {3pt}                           % 下方间距
  {\kaishu\upshape}               % 正文字体：中文楷体 + 英文直立
  {}                              % 缩进量
  {\bfseries\upshape}      % 头部字体（定理名）：中文楷体粗体 + 英文直立粗体
  {.}                             % 定理名后标点
  {0.5em}                         % 定理名后间距
  {}                              % 头部说明（留空）

% 应用自定义样式
\theoremstyle{kaiti-theorem}
\newtheorem{thm}{定理}[section]
\newtheorem{nota}[thm]{记号}
\newtheorem{question}{问题}
\newtheorem*{questionn}{思考}
\newtheorem{exe}{习题}[section]
\newtheorem{assump}[thm]{假设}
\newtheorem*{thmm}{定理}
\newtheorem{defn}{定义}[section]
\newtheorem{lem}[thm]{引理}
\newtheorem*{lemm}{引理}
\newtheorem{prop}[thm]{命题}
\newtheorem*{propp}{命题}
\newtheorem*{claim}{断言}
\newtheorem{cor}[thm]{推论}
\newtheorem{conj}[thm]{猜想}
\newtheorem{rmk}{注记}[section]
\newtheorem{ex}{例}[section]
\newtheorem{pf}{证明}
\newtheorem{problem}{作业题}
\numberwithin{equation}{section}





\DeclareMathOperator{\id}{id}
\DeclareMathOperator*{\esssup}{ess\,sup}
\renewcommand{\Re}{\textbf{Re}\,}  
\renewcommand{\Im}{\textbf{Im}\,}  
\newcommand{\R}{\mathbb{R}}  
\newcommand{\FH}{\mathbb{H}}  
\newcommand{\C}{\mathbb{C}}  
\newcommand{\Z}{\mathbb{Z}}
\newcommand{\X}{\mathbb{X}}
\newcommand{\N}{\mathbb{N}}
\newcommand{\Q}{\mathbb{Q}}
\newcommand{\T}{\mathbb{T}}
\newcommand{\TT}{\mathcal{T}}
\newcommand{\TP}{\overline{\partial}{}}
\newcommand{\TJ}{\langle\TP\rangle}
\newcommand{\TL}{\overline{\Delta}{}}
\newcommand{\curl}{\text{curl }}
\newcommand{\dive}{\text{div }}
\newcommand{\bd}[1]{\mathbf{#1}}  % for bolding symbols
\newcommand{\bds}[1]{\boldsymbol{#1}}  % for bolding symbols
\newcommand{\RR}{\mathcal{R}}      % for Real numbers
\newcommand{\sh}{\mathcal{S}}  
\newcommand{\sph}{\mathbb{S}}  
\newcommand{\Om}{\Omega}
\newcommand{\OmT}{{\Omega_T}}
\newcommand{\ut}{{U_T}}
\newcommand{\vp}{\varphi}
\newcommand{\Th}{\Theta}
\newcommand{\Lam}{\Lambda}
\newcommand{\Omb}{{\overline{\Omega}}}
\newcommand{\OmbT}{{\overline{\Omega_T}}}
\newcommand{\ub}{{\overline{U}}}
\newcommand{\ubt}{{\overline{U_T}}}
\newcommand{\q}{\quad}
\newcommand{\p}{\partial}
\newcommand{\pb}{\boldsymbol{\partial}}
\newcommand{\lap}{\Delta}
\newcommand{\dd}{\mathfrak{D}}
\newcommand{\DD}{\mathcal{D}}
\newcommand{\den}{{\bar\rho}_{0}}
\newcommand{\Dt}{{\p}_{t}+v^{k}{\p}_{k}}
\newcommand{\col}[1]{\left[\begin{matrix} #1 \end{matrix} \right]}
\newcommand{\noi}{\noindent}
\newcommand{\nab}{\nabla}
\newcommand{\cnab}{\overline{\nab}}
\newcommand{\cp}{\overline{\partial}{}}
\newcommand{\dx}{\,\mathrm{d}x}
\newcommand{\dxx}{\,\mathrm{d}\bds{x}}
\newcommand{\xxi}{\boldsymbol{\xi}}
\newcommand{\eeta}{\boldsymbol{\eta}}
\newcommand{\dxxi}{\,\mathrm{d}\boldsymbol{\xi}}
\newcommand{\deeta}{\,\mathrm{d}\boldsymbol{\eta}}
\newcommand{\dxi}{\,\mathrm{d}\xi}
\newcommand{\deta}{\,\mathrm{d}\eta}
\newcommand{\dy}{\,\mathrm{d}y}
\newcommand{\dyy}{\,\mathrm{d}\bds{y}}
\newcommand{\dz}{\,\mathrm{d}z}
\newcommand{\dzz}{\,\mathrm{d}\bds{z}}
\newcommand{\dph}{\,\mathrm{d}\varphi}
\newcommand{\dt}{\,\mathrm{d}t}
\newcommand{\dr}{\,\mathrm{d}r}
\newcommand{\dtt}{\,\mathrm{d}\tau}
\newcommand{\dS}{\,\mathrm{d}S}
\newcommand{\dss}{\,\mathrm{d}\sigma}
\newcommand{\dmu}{\,\mathrm{d}\mu}
\newcommand{\dnu}{\,\mathrm{d}\nu}
\newcommand{\dV}{\,\mathrm{d}V}
\newcommand{\ds}{\,\mathrm{d}s}
\newcommand{\drr}{\,\mathrm{d}\rho}
\newcommand{\dist}{\text{dist }}
\newcommand{\vol}{\text{vol}\,}
\newcommand{\volo}{\text{vol}(\Omega)}
\newcommand{\spt}{\text{Spt}\,}
\newcommand{\Tr}{\text{Tr}\,}
\newcommand{\lee}{\langle}
\newcommand{\ree}{\rangle}
\newcommand{\xee}{\langle\xi\rangle}
\newcommand{\xxee}{\langle\boldsymbol{\xi}\rangle}
\newcommand{\xeta}{\langle\eta\rangle}
\newcommand{\xeeta}{\langle\boldsymbol{\eta}\rangle}
\newcommand{\kk}{\kappa}
\newcommand{\lam}{\lambda}
\newcommand{\io}{\int_{\Omega}}
\newcommand{\iu}{\int_{U}}
\newcommand{\is}{\int_{\Sigma}}
\newcommand{\ipo}{\int_{\partial\Omega}}
\newcommand{\ipu}{\int_{\partial U}}
\newcommand{\ig}{\int_{\Gamma}}
\newcommand{\ir}{\int_{\R}}
\newcommand{\ird}{\int_{\R^d}}
\newcommand{\irn}{\int_{\R^n}}
\newcommand{\fpi}{\frac{1}{(2\pi)^{\frac{d}{2}}}}
\newcommand{\ddt}{\frac{\mathrm{d}}{\mathrm{d}t}}
\newcommand{\ddx}{\frac{\mathrm{d}}{\mathrm{d}x}}
\newcommand{\eps}{\varepsilon}
\providecommand{\jump}[1]{\left\llbracket #1 \right\rrbracket }
\providecommand{\len}[1]{\lee #1 \ree }
\providecommand{\ino}[1]{\left\| #1 \right\| }
\providecommand{\bno}[1]{\left| #1 \right| }
\newcommand{\red}{\textcolor{red}}
\newcommand{\green}{\textcolor{green}}
\newcommand{\blue}{\textcolor{blue}}
\newcommand{\nnr}{{[n+1]}}
\newcommand{\nnn}{{[n]}}
\newcommand{\nnl}{{[n-1]}}
\newcommand{\nnll}{{[n-2]}}	
\newcommand{\mmr}{{[m+1]}}
\newcommand{\mmm}{{[m]}}
\newcommand{\mml}{{[m-1]}}
\newcommand{\mmll}{{[m-2]}}

%---------------Special variables--------------
\newcommand{\CC}{\mathfrak{C}}  
\newcommand{\NN}{\mathbf{N}}
\newcommand{\LH}{\mathcal{L}}
\newcommand{\HH}{\mathcal{H}}  
\newcommand{\EE}{\mathcal{E}}  
\newcommand{\FT}{\mathcal{F}}  
\newcommand{\fT}{\mathbf{T}}  
\newcommand{\FA}{\mathbf{A}}  
\newcommand{\MH}{\mathcal{M}}
\newcommand{\NH}{\mathcal{N}}
\newcommand{\PH}{\mathcal{P}}
\newcommand{\VH}{\mathcal{V}}
\newcommand{\XX}{\mathfrak{X}}
\newcommand{\xx}{{\bds{x}}}
\newcommand{\bx}{\mathbf{x}}
\newcommand{\by}{\mathbf{y}}
\newcommand{\bp}{\mathbf{p}}
\newcommand{\bq}{\mathbf{q}}
\newcommand{\pp}{{\bds{p}}}
\newcommand{\qq}{{\bds{q}}}
\newcommand{\uu}{\mathbf{u}}
\newcommand{\ww}{\mathbf{w}}
\newcommand{\vv}{\mathbf{v}}
\newcommand{\yy}{{\bds{y}}}
\newcommand{\zz}{{\bds{z}}}
\newcommand{\zo}{\mathbf{0}}
\newcommand{\ee}{\mathbf{e}}
\newcommand{\fa}{\mathbf{a}}
\newcommand{\fb}{\mathbf{b}}
\newcommand{\fc}{\mathbf{c}}
\newcommand{\ff}{\bds{f}}
\newcommand{\fr}{\mathbf{r}}
\newcommand{\fh}{\mathbf{h}}
\newcommand{\fm}{\mathbf{m}}
\newcommand{\fg}{\mathbf{g}}
\newcommand{\FF}{\mathbf{F}}
\newcommand{\GG}{\mathbf{G}}
\newcommand{\BB}{\mathbf{B}}



\begin{document}
\title{\LARGE{\bf  2026年春季学期现代偏微分方程作业五}}
\author{波动方程}
\date{不用提交}

\maketitle

\begin{problem}[习题7.1.2及后续]
设$D\in\R$为常数，$U\subset\R^d$ 是一个具有光滑边界的区域，且 $\vp ,\psi\in C_c^\infty(U)$.
\begin{equation*}
\begin{cases}
\p_t^2 u+D\p_t u-\lap u=0~~~&\text{ in }(0,T)\times U,\\
u=\varphi,~\p_t u=\psi~~~&\text{ on }\{t=0\}\times U,\\
u=0~~~&\text{ on }[0,T]\times\p U.
\end{cases}
\end{equation*}
\begin{itemize}
\item [(1)] 证明：对于\textbf{任意}常数 $D\in\R$，如上方程至多有一个光滑解 $u\in C^{\infty}([0,T)\times \ub)$.
\item [(2)] 设$D=1$，且已知方程的整体光滑解存在（即可以延拓到$t=+\infty$），定义能量 $E_0(t):=\frac12\int_U (\p_tu(t,\xx))^2+|\nab u(t,\xx)|^2\dxx$. 证明：$\lim\limits_{t\to+\infty} E_0(t)=0$.
\end{itemize}

提示：(2) 对 $\eps>0$，先考虑 $E_\eps(t):=E_0(t)+\eps\iu u\p_t u\dxx$ 的能量估计，再证明 $\eps$ 适当小时 $E_\eps, E_0$ 是可比较的。
\end{problem}

\begin{problem}[习题7.2.2] 
考虑半线性波动方程的初值问题
\begin{equation}
\begin{cases}
\p_t^2 u -\lap  u +f( u )=0&~~t>0,~\xx\in\R^d;\\
 u (0,\xx)= u_0(\xx),~\p_t  u (0,x)=u_1(\xx)&~~\xx\in\R^d.
\end{cases}
\end{equation} 这里 $f$ 是连续函数，且假设 $ u $ 在 $|\xx|\to\infty$ 时趋于零。
\begin{itemize}
\item [(1)] 证明如下$E(t)$关于时间是守恒量
\[
E(t):=\int_{\R^d}\frac12(|\p_t  u |^2+|\nab  u |^2) + F( u )\dxx,\q\q F(u):=\int_0^u f(s)\ds.
\] 
\item [(2)] 给定 $t_0>0$ 和 $\xx_0\in\R^d$ 并定义带有顶点 $(\xx_0,t_0)$ 的时间倒向光锥 (time-backward light cone) 为
\[
K(\xx_0,t_0):=\{(t,\xx)\in[0,\infty)\times \R^d : 0\leqslant t\leqslant t_0,~|\xx-\xx_0|\leqslant t_0-t\},
\]  $K(\xx_0,t_0)$ 边界的弯曲部分为
\[
\Gamma(\xx_0,t_0):=\{(t,\xx)\in  [0,\infty)\times\R^d: 0\leqslant t\leqslant t_0,~|\xx-\xx_0|= t_0-t\}.
\] 定义光锥上的能量通量(energy flux)为
\[
e(t):=\int_{B(\xx_0,t_0-t)}\frac12(|\p_t  u |^2+|\nab  u |^2) +F( u )\dxx~~(0\leqslant t\leqslant t_0).
\] 证明：
\begin{equation}
\frac{1}{\sqrt{2}}\int_{\Gamma(\xx_0,t_0)}\frac12|(\p_t  u )\nu-\nab  u |^2+F( u )\dS=e(0)
\end{equation}其中 $\nu:=\frac{\xx-\xx_0}{|\xx-\xx_0|}$.
\item [(3)] 设 $F\geqslant 0$，利用 (2) 证明如果$ u_0,  u_1$在$B(\xx_0,t_0)$中恒为零，那么$ u $在光锥$K(\xx_0,t_0)$ 内恒为零。
\item [(4)] 证明 (3) 对于拟线性波动方程 $\p_t^2  u  -\lap  u  + f( u ,\pb  u )=0$ 的光滑解 $ u $ 也成立，其中 $f$ 是局部 Lipschitz 连续的且 $f(0,\mathbf{0})=0$. 
\end{itemize}

提示：(2) 计算 $e'(t)$ 并将其与待证等式的左端进行比较，尤其要思考系数 $1/\sqrt{2}$ 是如何出现的。(4) 考虑 $E(t)=\int_{B(\xx_0,t_0-t)}\frac12(|\p_t  u |^2+|\nab  u |^2) +  u^2\dxx$ 并注意到存在依赖于 $\| u ,\pb  u \|_{L^\infty}$ 的常数 $C>0$ 使得 $|f( u ,\pb u )|\leqslant C(| u |+|\pb u |)$成立。
\end{problem}


\begin{problem}[习题7.3.2]
设$\bds{\phi}:\R\times\R^d\to\sph^m\subset \R^{m+1}$满足如下波映射(wave map)方程
\[
\begin{cases} 
\square \bds{\phi}=  \bds{\phi}\left(\p_t  \bds{\phi}^\top\cdot \p_t  \bds{\phi} - \sum\limits_{i=1}^d \p_i \bds{\phi}^\top \cdot \p_i  \bds{\phi} \right), & (t, \xx) \in (0, T) \times \R^d, \\ 
 \bds{\phi}(0, \xx) =  \bds{\phi}_0(\xx), \quad \p_t  \bds{\phi}(0, \xx) =  \bds{\phi}_1(\xx). 
\end{cases}
\]初值满足$| \bds{\phi}_0 |^2=1,  \bds{\phi}_1^\top \cdot  \bds{\phi}_0 = 0$，并假设具有紧支集。
\begin{itemize}
\item [(1)]  证明如下能量 $E(t)$ 是守恒量$$E(t):=\frac12\ird \left(|\p_t  \bds{\phi}(t,\xx)|^2 + \sum_{i=1}^d |\p_i  \bds{\phi}(t,\xx)|^2 \dxx\right).$$
\item [(2)]  设$d=1,m=2$，并假设光滑初值满足：存在$\yy\in \sph^2$ 使得 $( \bds{\phi}_0-\yy, \bds{\phi}_1)\in H^2(\R)\times H^1(\R)$. 证明：此时波映射方程有光滑的整体解。这里可以默认爆破准则(7.3.18)成立，即
\[
\text{解的极大存在时间}T_*<\infty\implies\limsup_{t\to T_*}\sum\limits_{|\alpha|\leqslant 1} \|\pb^\alpha  \bds{\phi}(t,\cdot)\|_{L^\infty(\R^d)}=\infty.
\]
\end{itemize}
\end{problem}


\begin{problem}[问题7.4.1 (Levine凹性方法)]
考虑如下聚焦次临界波方程
\[
\begin{cases} 
\p_t^2 u - \lap u = u^3, & (t, \xx) \in [0, T) \times \R^3, \\ 
u(0, \xx) = u_0(\xx), \quad \p_t u(0, \xx) = u_1(\xx)&\xx\in\R^3. \end{cases}
\]其中初值$u_0,u_1\in C_c^\infty(\R^3)$支于球$B(0,R)$内，$R>0$给定，且满足$E(0):=\int_{\R^3}\frac12(|\nab u_0|^2+|u_1|^2)-\frac14 |u_0|^4 \dxx<0$.
\begin{itemize}
\item [(1)] 证明：存在常数 $A>0$ 使得 $\|\nabla u_0\|_{L^2}>A$，因此满足 $E(0)<0$ 的初值不可能是小初值。
\item [(2)] 令 $I(t):=\int_{\R^3} u^2\dxx$，证明：$I''(t)\geqslant 6\int_{\R^3}(\p_t u)^2\dxx-8E(0)$，进而存在 $t_1\geqslant 0$ 使得对任意 $t\geqslant t_1$ 有 $I'(t)>0$.
\item [(3)] 令 $J(t):=I(t)^{-1/2}$，证明: $J''(t)<0$, 且对任意 $t\geqslant t_1$ 有 $J'(t)<0$. 据此得出存在 $T_*<\infty$ 使得 $J(T_*)=0$, 即 $I(T_*)=+\infty$.
\end{itemize}
\end{problem}




\end{document}
